% \documentclass[conference]{IEEEtran}
\documentclass{article}
\usepackage[T1]{fontenc}
\usepackage[utf8]{inputenc}
\usepackage{listings}
\usepackage{fontspec}
\setmonofont{DejaVu Sans Mono}
\usepackage{amssymb}
\usepackage{amsmath}
\usepackage[letterpaper, margin=1in]{geometry}

\usepackage{color}
\definecolor{keywordcolor}{rgb}{0.7, 0.1, 0.1}   % red
\definecolor{tacticcolor}{rgb}{0.0, 0.1, 0.6}    % blue
\definecolor{commentcolor}{rgb}{0.4, 0.4, 0.4}   % grey
\definecolor{symbolcolor}{rgb}{0.0, 0.1, 0.6}    % blue
\definecolor{sortcolor}{rgb}{0.1, 0.5, 0.1}      % green
\definecolor{attributecolor}{rgb}{0.7, 0.1, 0.1} % red

\DeclareMathOperator{\spann}{span}
\newcommand\spanset[1]{\ensuremath\spann(#1)}

\def\lstlanguagefiles{lstlean.tex}
% set default language
\lstset{language=lean}
\usepackage{graphicx} % Required for inserting images

\title{L$\exists \forall$N Formalization: Spans are Subspaces}
\author{Grant Yang}
\date{December 2025}

\begin{document}

\maketitle

We would like to formalize the idea that spans are subspaces in Lean. We start by specifying that there is some type \lstinline@V@ of vectors that forms a module over the reals. We also need to specify that \lstinline@V@ has an additive commutative group structure because Mathlib's definition of \lstinline@Module@ depends on it.

\begin{lstlisting}
variable {V} [AddCommGroup V] [Module ℝ V]
\end{lstlisting}

In Lean, sets are simply functions that return true on elements in the set and false on elements not in the set. Actually, these functions do not really return a boolean true or false value, but instead return something called a proposition, \lstinline@Prop@, which is much more powerful. For example, we can define the set of zero and the even natural numbers like so:

\begin{lstlisting}
def evens_union_zero : Set ℕ := fun x ↦ ∃(n : ℕ), 2 * n = x
\end{lstlisting}

Notice how \lstinline@Set N@ is synonymous with \lstinline@N -> Prop@, and our code corresponds exactly to the expression $\{x\in \mathbb{N} \mid \exists(n\in \mathbb N), 2n=x\}$ in set builder notation. 

For any $S \subseteq V$, we can now define $\spann S$ to be the set of vectors that are linear combinations of vectors in $S$. A vector $\vec{x}$ is a linear combination of the vectors in $S$ if there exist both a list $a_1,a_2,\dots,a_n$ of real numbers and a list $\vec{v_1},\vec{v_2},\dots,\vec{v_n}$ of vectors in $S$ such that $\vec{x}=a_1 \vec{v_1} + a_2 \vec{v_2}+\dots+ a_n \vec{v_n}$.

For our purposes, we can implement this definition using linked lists of real numbers and vectors. We can zip these lists together into a single list of $\mathbb{R} \times S$, but due to technical details it is somewhat clunky to express the type of elements of the set $S$. Instead, it is easier to use $C \subseteq \mathbb{R}\times V$ and require that we have $\forall (c \in C), \pi_2(c)\in S$.

\begin{lstlisting}

def span (S : Set V) : Set V :=
  {x : V | ∃(C : List (ℝ×V)),
    (sumUp C = x) ∧ (allin C S) }


-- proposition that all vectors (v : V) in the list are also elements of the set S.
def allin {V} (ls : List (ℝ × V)) (S : Set V) : Prop := match ls with
| List.cons ⟨_, v⟩ rest => v ∈ S ∧ allin rest S
| List.nil => True

def sumUp (vecs : List (ℝ × V)) : V := match vecs with
  | List.cons ⟨c, v⟩ rest => c • v + sumUp rest
  | List.nil => 0

\end{lstlisting}

Immediately from this definition, we can already prove that the zero vector is in all spans since we have defined \lstinline@sumUp [] = 0@, so we can always prove ``there exists'' a linear combination for $0$ by providing the empty list as an example. In the following code, the expression appearing after \lstinline@|-@ is where I show the current tactic state, and it is the proposition we are trying to prove.

\begin{lstlisting}
theorem zero_in_span : ∀(S : Set V), 0 ∈ span S := by
  intro S
  unfold span
  let ex : List (ℝ × V) := List.nil
  dsimp
  -- ⊢ ∃ C, sumUp C = 0 ∧ allin C S
  fconstructor -- constructor for ∃
  · exact ex -- example empty list
  · -- unfold definitions to show it's a linear combo for 0
    unfold ex sumUp allin
    simp
\end{lstlisting}

Note that because ``for all'' statements correspond to function types (by the Curry-Howard correspondence), the \lstinline@intro@ tactic allows us to prove a ``for all'' statement by introducing an argument name and constructing the function explicitly.

However, in order to use this definition of span, we must first prove some basic theorems about \lstinline@allin@ and \lstinline@sumUp@ as they are custom operations that we have defined. Let us start by showing that they distribute over list concatenation. If we have two lists \lstinline@la lb : List (ℝ × V)@, then we have \lstinline@allin (la ++ lb) S@ if and only if both \lstinline@allin la S@ and \lstinline@allin lb S@. Similarly, \lstinline@(sumUp la) + (sumUp lb) = sumUp (la ++ lb)@.

Let us start with the proof for the \lstinline@sumUp@ property because it is simpler. Because we are using linked lists, this can be proved inductively.

\begin{lstlisting}

theorem sumup_concat : ∀(a b : List (ℝ×V)),
    sumUp (a++b) = sumUp a + sumUp b := by
  intros a b
  -- linked lists are an inductive type so we can do induction on a
  induction a with
  | nil => 
    -- base case: ⊢ sumUp ([] ++ b) = sumUp [] + sumUp b
    simp [sumUp] -- just unfold the definition of sumUp!
  | cons head tail ih =>
    -- inductive hypothesis ih : sumUp (tail ++ b) = sumUp tail + sumUp b
    -- inductive case: ⊢ sumUp (head :: tail ++ b) = sumUp (head :: tail) + sumUp b
    simp [sumUp]
    -- ⊢ head.1 • head.2 + sumUp (tail ++ b) = head.1 • head.2 + sumUp tail + sumUp b
    rw [ih]
    rw [add_assoc]
\end{lstlisting}

The proof for the analogous property on \lstinline@allin@ follows basically the same structure, but is a bit more involved:


\begin{lstlisting}
theorem allin_concat {V} : ∀(la lb : List (ℝ × V)) (S : Set V),
    allin la S ∧ allin lb S ↔ allin (la ++ lb) S := by
  intros la lb S
  induction la with
  | nil =>
    -- ⊢ allin [] S ∧ allin lb S ↔ allin ([] ++ lb) S
    simp [allin] -- just unfold the definition of allin
  | cons head tail ih =>
    -- ih : allin tail S ∧ allin lb S ↔ allin (tail ++ lb) S
    -- ⊢ allin (head::tail) S ∧ allin lb S ↔ allin (head::tail ++ lb) S
    simp [allin]
    rw [and_assoc]
    -- ⊢ head.2 ∈ S ∧ allin tail S ∧ allin lb S ↔ head.2 ∈ S ∧ allin (tail ++ lb) S
    constructor -- construct the Iff
    · -- forward direction of Iff
      intros m -- head.2 ∈ S ∧ allin tail S ∧ allin lb S
      constructor -- And constructor
      · exact m.left -- head.2 ∈ S
      · exact ih.mp m.right -- allin (tail ++ lb) S
    · -- backward direction of Iff
      intros m -- head.2 ∈ S ∧ allin (tail ++ lb) S
      constructor -- And constructor
      · exact m.left -- head.2 ∈ S
      · exact ih.mpr m.right -- allin tail S ∧ allin lb S
\end{lstlisting}

With these two theorems, we are now able to prove that spans are additively closed! To state this theorem, we use a \lstinline@Subtype@ to express ``elements of the set S'' as a type. Instances of the subtype \lstinline@{v : V // v \in span S}@ are simply an object \lstinline@v : V@ together with a proof that \lstinline@v \in span S@. To prove that spans are additively closed, we show that $\forall(\vec a, \vec b\in \spann S), \vec a + \vec b \in \spann S$. We can use our definition of span to choose linear combinations for $\vec a$ and $\vec b$ and simply concatenate them together to make a linear combination for $\vec{a} + \vec{b}$, using the two theorems we just proved to show that our concatenated linear combination sums up to $\vec{a} + \vec{b}$ and still only uses vectors from $S$.

\begin{lstlisting}
theorem span_add_closed : ∀(S : Set V) (a b : {v : V // v ∈ span S}),
    a.val + b.val ∈ span S := by
  intros S a b
  let pa := a.prop -- a ∈ span S
  let pb := b.prop -- b ∈ span S
  unfold span at pa; dsimp at pa
  unfold span at pb; dsimp at pb

  -- pa : ∃ C, sumUp C = ↑a ∧ allin C S
  --      ⟨xa, ⟨sumUp xa = ↑a, allin xa S⟩⟩
  rcases pa with ⟨xa, ⟨ha_sum, ha_sub⟩⟩
  -- pb : ∃ C, sumUp C = ↑b ∧ allin C S
  --      ⟨xb, ⟨sumUp xb = ↑b, allin xb S⟩⟩
  rcases pb with ⟨xb, ⟨hb_sum, hb_sub⟩⟩
  have sumeq: sumUp (xa ++ xb) = a.val + b.val := by
    calc
      _ = sumUp xa + sumUp xb := by exact sumup_concat xa xb
      _ = a.val + b.val := by rw [ha_sum, hb_sum]
  have ain: allin (xa ++ xb) S := by
    exact (allin_concat xa xb S).mp ⟨ha_sub, hb_sub⟩

  unfold span
  dsimp
  -- ⊢ ∃ C, sumUp C = ↑a + ↑b ∧ allin C S
  constructor
  · exact ⟨sumeq, ain⟩
\end{lstlisting}

Now all that is left is to show that spans are multiplicatively closed. To do this, we can show that for all $\vec{u} \in \spann S$ and $r \in \mathbb R$, we can obtain a linear combination for $r \cdot \vec u$ by multiplying the real part of the linear combination for $\vec{u}$ by $r$. We define a new operation $\text{lsmul} : \mathbb{R} \to \text{List}\,(\mathbb{R}\times V )\to \text{List}\,(\mathbb{R} \times V)$ that performs this scalar multiplication of linear combinations, and we need to show that if $C : \text{List}\,(\mathbb{R} \times V)$ is some linear combination for $\vec{u}$, then \lstinline@lsmul r C@ is a linear combination for $r \cdot \vec{u}$.

\begin{lstlisting}
def lsmul (r : ℝ) (ls : List (ℝ × V)) : List (ℝ × V) :=
    List.map (fun ⟨a, b⟩ ↦ ⟨r • a, b⟩) ls

-- dedicated theorem just for rewriting between lsmul and its expanded definition. Is there a better way?
theorem lsmul_def : ∀(ls : List (ℝ×V)) (r : ℝ), lsmul r ls =
    List.map (fun ⟨a, b⟩ ↦ ⟨r * a, b⟩) ls := by simp [lsmul]
\end{lstlisting}

Now that we have defined a custom operation, we once again need to prove a few theorems about it. First, let us show that \lstinline@sumUp@ and \lstinline@lsmul@ obey the distributive property as expected, such that \lstinline@sumUp (lsmul r ls) = r • (sumUp ls)@.

\begin{lstlisting}
-- helper theorem because real multiplication on ℝ (*) and scalar multiplication on modules (•) use different notation.
theorem real_mul_smul : ∀(a b : ℝ), a * b = a • b := by simp

theorem sumup_lin : ∀(ls : List (ℝ×V)) (r : ℝ),
    sumUp (lsmul r ls) = r • (sumUp ls) := by
  intro ls r
  induction ls with
  | nil =>
    unfold lsmul sumUp
    simp -- empty list: sumUp [] = 0, and 0 = x • 0
  | cons head tail ih =>
    -- ih : sumUp (lsmul r tail) = r • sumUp tail
    unfold sumUp lsmul; dsimp
    -- ⊢ (r * head.1) • head.2 + sumUp (List.map (fun x ↦ (r * x.1, x.2)) tail) = r • (head.1 • head.2 + sumUp tail)
    rw [← lsmul_def, ih]
    -- ⊢ (r * head.1) • head.2 + r • sumUp tail = r • (head.1 • head.2 + sumUp tail)
    rw [smul_add, real_mul_smul, smul_assoc]
\end{lstlisting}

Finally, we must prove that \lstinline@lsmul@ does not modify the vector part of the pairs, so we are still only using vectors from $S$. We are actually able to prove this statement as an if and only if:

\begin{lstlisting}
theorem lsmul_allin : ∀(ls : List (ℝ×V)) (S: Set V) (r : ℝ),
    allin ls S ↔ allin (lsmul r ls) S := by
  intros ls S r
  constructor -- Iff
  · induction ls with -- forward direction of Iff
    | nil => -- ⊢ allin [] S → allin (lsmul r []) S
      unfold lsmul List.map; simp
    | cons head tail ih =>
      -- ih : allin tail S → allin (lsmul r tail) S
      -- ⊢ allin (head :: tail) S → allin (lsmul r (head :: tail)) S
      unfold allin lsmul; dsimp
      rw [← lsmul_def]
      -- ⊢ head.2 ∈ S ∧ allin tail S → head.2 ∈ S ∧ allin (lsmul r tail) S
      intro ⟨elem, rest⟩
      exact ⟨elem, ih rest⟩
  · induction ls with -- backward direction of Iff
    | nil => -- ⊢ allin (lsmul r []) S → allin [] S
      unfold lsmul List.map; simp
    | cons head tail ih =>
      -- ih : allin (lsmul r tail) S → allin tail S
      -- ⊢ allin (lsmul r (head :: tail)) S → allin (head :: tail) S
      unfold allin lsmul; dsimp
      rw [← lsmul_def]
      -- ⊢ head.2 ∈ S ∧ allin (lsmul r tail) S → head.2 ∈ S ∧ allin tail S
      intro ⟨elem, rest⟩
      exact ⟨elem, ih rest⟩
\end{lstlisting}

Now, we can prove that spans are multiplicatively closed.

% \newpage{}
\begin{lstlisting}
theorem span_mul_closed : ∀(S : Set V) (r : ℝ)
    (u : {v : V // v ∈ span S}), r • u.val ∈ span S := by
  intros S r u
  unfold span; dsimp
  -- ⊢ ∃ C, sumUp C = r • ↑u ∧ allin C S

  let pu := u.prop
  unfold span at pu; dsimp at pu
  -- pu : ∃ C, sumUp C = ↑u ∧ allin C S 
  --       ⟨lu, sumUp lu = ↑u ∧ allin lu S⟩
  rcases pu with ⟨lu, pu⟩
  let ex: List (ℝ × V) := lsmul r lu
  have sum_val : sumUp ex = r • u := by
    unfold ex
    -- ⊢ sumUp (lsmul r lu) = r • ↑u
    rw [sumup_lin, pu.left]
  have allin_proof : allin ex S := by
    unfold ex
    -- ⊢ allin (lsmul r lu) S
    exact (lsmul_allin lu S r).mp pu.right
  constructor -- for ∃
  · exact ⟨sum_val, allin_proof⟩
\end{lstlisting}

This mostly completes the formalization! All that is left is to use these theorems to formally construct a \lstinline@Module@ instance on the \lstinline@{v : V // v \in span S}@ subtype. We have to do some ceremony to bridge the gap between the \lstinline@Subtype@ and the type of vectors \lstinline@V@, but completing the construction is mostly straightforward.

\begin{lstlisting}
-- define the zero vector on the span subspace
instance (S : Set V) : Zero {v : V // v ∈ span S} where
  zero := Subtype.mk 0 <| zero_in_span S

-- vector negation
instance (S : Set V) : Neg {v : V // v ∈ span S} where
  neg (a : {v : V // v ∈ span S}) :=
    Subtype.mk (-a.val) <| by
      have sub : -a.val = (-1 : ℝ) • a.val := by simp
      rw [sub]
      exact span_mul_closed S (-1) a

-- vector addition
instance (S : Set V) : Add {v : V // v ∈ span S} where
  add (a b : {v: V // v ∈ span S}) :=
    Subtype.mk (a.val + b.val) <| span_add_closed S a b

@[simp]
theorem subtype_ext_add : ∀(S: Set V) (a b : {v : V // v ∈ span S}),
    (a + b).val = a.val + b.val := by intros S a b; rfl

@[simp]
theorem subtype_ext_neg : ∀(S : Set V) (a : {v : V // v ∈ span S}),
  (-a).val = -(a.val) := by intros S a; rfl

@[simp]
theorem subtype_ext_zero : ∀(S : Set V),
  (0 : {v : V // v ∈ span S}).val = 0 := by intro S; rfl

-- vector addition is an abelian group
instance (S : Set V) : AddCommGroup {v: V // v ∈ span S} where
  add_assoc := by intro a b c; ext; simp; rw [add_assoc]
  zero_add := by intro z; ext; simp;
  add_zero := by intro a; ext; simp;
  add_comm := by intro a b; ext; simp; rw [add_comm]
  neg_add_cancel := by intro a; ext; simp;
  nsmul := nsmulRec
  zsmul := zsmulRec

-- scalar multiplication
instance (S : Set V) : SMul ℝ {v : V // v ∈ span S} where
  smul (r : ℝ) (u : {v : V // v ∈ span S}) :=
    Subtype.mk (r • u.val) <| span_mul_closed S r u

@[simp]
theorem subtype_ext_smul : ∀(S: Set V) (r: ℝ)
    (u : {v : V // v ∈ span S}), (r • u).val = r • u.val := by
  intros S a b; rfl

-- we have the module axioms: spans are subspaces!
instance (S : Set V) : Module ℝ {v : V // v ∈ span S} where
  one_smul := by intro b; ext; simp
  mul_smul := by intros a b c; ext; simp; rw [mul_smul]
  smul_zero := by intros a; ext; simp;
  smul_add := by intros a b c; ext; simp
  add_smul := by intros a b c; ext; simp; rw [add_smul]
  zero_smul := by intros a; ext; simp
\end{lstlisting}

QED! $\square$


\onecolumn
\newpage{}
\section{Appendix}
    \begin{lstlisting}
import Mathlib.Algebra.Module.Basic
import Mathlib.Data.Real.Basic
import Mathlib.Tactic

-- proposition that all vectors (v : V) in the list are also elements of the set S.
def allin {V} (ls : List (ℝ × V)) (S : Set V) : Prop := match ls with
| List.cons ⟨_, v⟩ rest => v ∈ S ∧ allin rest S
| List.nil => True

theorem allin_concat {V} : ∀(la lb : List (ℝ × V)) (S : Set V),
    allin la S ∧ allin lb S ↔ allin (la ++ lb) S := by
  intros la lb S
  induction la with
  | nil =>
    -- ⊢ allin [] S ∧ allin lb S ↔ allin ([] ++ lb) S
    simp [allin] -- just unfold the definition of allin
  | cons head tail ih =>
    -- ih : allin tail S ∧ allin lb S ↔ allin (tail ++ lb) S
    -- ⊢ allin (head::tail) S ∧ allin lb S ↔ allin (head::tail ++ lb) S
    simp [allin]
    rw [and_assoc]
    -- ⊢ head.2 ∈ S ∧ allin tail S ∧ allin lb S ↔ head.2 ∈ S ∧ allin (tail ++ lb) S
    constructor -- construct the Iff
    · -- forward direction of Iff
      intros m -- head.2 ∈ S ∧ allin tail S ∧ allin lb S
      constructor -- And constructor
      · exact m.left -- head.2 ∈ S
      · exact ih.mp m.right -- allin (tail ++ lb) S
    · -- backward direction of Iff
      intros m -- head.2 ∈ S ∧ allin (tail ++ lb) S
      constructor -- And constructor
      · exact m.left -- head.2 ∈ S
      · exact ih.mpr m.right -- allin tail S ∧ allin lb S

/- this is not used in the "subspace closed proof" but
  i wrote it as i was trying to figure out how to do inductive proofs -/
theorem allin_ordering {V} : ∀(ls : List (ℝ × V)) (A B : Set V),
    A ⊆ B → allin ls A → allin ls B := by
  intros ls A B subset allA
  induction ls with
  | nil => simp [allin]
  | cons head rest ih =>
    simp [allin] at allA
    simp [allin]
    constructor
    · apply subset
      exact allA.left
    · apply ih
      exact allA.right

-- helper theorem because real multiplication on ℝ (*)
-- and scalar multiplication on modules (•) use different notation.
theorem real_mul_smul : ∀(a b : ℝ), a * b = a • b := by simp

namespace Vect

variable {V} [AddCommGroup V] [Module ℝ V]

def sumUp (vecs : List (ℝ × V)) : V := match vecs with
  | List.cons ⟨c, v⟩ rest => c • v + sumUp rest
  | List.nil => 0

def span (S : Set V) : Set V :=
  {x : V | ∃(C : List (ℝ×V)),
    (sumUp C = x) ∧ (allin C S) }

theorem sumup_concat : ∀(a b : List (ℝ×V)),
    sumUp (a++b) = sumUp a + sumUp b := by
  intros a b
  -- linked lists are an inductive type so we can do induction on a
  induction a with
  | nil =>
    -- base case: ⊢ sumUp ([] ++ b) = sumUp [] + sumUp b
    simp [sumUp] -- just unfold the definition of sumUp!
  | cons head tail ih =>
    -- inductive hypothesis ih : sumUp (tail ++ b) = sumUp tail + sumUp b
    -- inductive case: ⊢ sumUp (head :: tail ++ b) = sumUp (head :: tail) + sumUp b
    simp [sumUp]
    -- ⊢ head.1 • head.2 + sumUp (tail ++ b) = head.1 • head.2 + sumUp tail + sumUp b
    rw [ih]
    rw [add_assoc]

def lsmul (r : ℝ) (ls : List (ℝ × V)) : List (ℝ × V) :=
    List.map (fun ⟨a, b⟩ ↦ ⟨r • a, b⟩) ls

-- dedicated theorem just for rewriting between lsmul and its expanded definition.
-- Is there a better way?
omit [AddCommGroup V] [Module ℝ V] in
theorem lsmul_def : ∀(ls : List (ℝ×V)) (r : ℝ), lsmul r ls =
    List.map (fun ⟨a, b⟩ ↦ ⟨r * a, b⟩) ls := by simp [lsmul]

theorem sumup_lin : ∀(ls : List (ℝ×V)) (r : ℝ),
    sumUp (lsmul r ls) = r • (sumUp ls) := by
  intro ls r
  induction ls with
  | nil =>
    unfold lsmul sumUp
    simp -- empty list: sumUp [] = 0, and 0 = x • 0
  | cons head tail ih =>
    -- ih : sumUp (lsmul r tail) = r • sumUp tail
    unfold sumUp lsmul; dsimp
    -- ⊢ (r * head.1) • head.2 + sumUp (List.map (fun x ↦ (r * x.1, x.2)) tail) = r • (head.1 • head.2 + sumUp tail)
    rw [← lsmul_def, ih]
    -- ⊢ (r * head.1) • head.2 + r • sumUp tail = r • (head.1 • head.2 + sumUp tail)
    rw [smul_add, real_mul_smul, smul_assoc]

omit [AddCommGroup V] [Module ℝ V] in
theorem lsmul_allin : ∀(ls : List (ℝ×V)) (S: Set V) (r : ℝ),
    allin ls S ↔ allin (lsmul r ls) S := by
  intros ls S r
  constructor -- Iff
  · induction ls with -- forward direction of Iff
    | nil => -- ⊢ allin [] S → allin (lsmul r []) S
      unfold lsmul List.map; simp
    | cons head tail ih =>
      -- ih : allin tail S → allin (lsmul r tail) S
      -- ⊢ allin (head :: tail) S → allin (lsmul r (head :: tail)) S
      unfold allin lsmul; dsimp
      rw [← lsmul_def]
      -- ⊢ head.2 ∈ S ∧ allin tail S → head.2 ∈ S ∧ allin (lsmul r tail) S
      intro ⟨elem, rest⟩
      exact ⟨elem, ih rest⟩
  · induction ls with -- backward direction of Iff
    | nil => -- ⊢ allin (lsmul r []) S → allin [] S
      unfold lsmul List.map; simp
    | cons head tail ih =>
      -- ih : allin (lsmul r tail) S → allin tail S
      -- ⊢ allin (lsmul r (head :: tail)) S → allin (head :: tail) S
      unfold allin lsmul; dsimp
      rw [← lsmul_def]
      -- ⊢ head.2 ∈ S ∧ allin (lsmul r tail) S → head.2 ∈ S ∧ allin tail S
      intro ⟨elem, rest⟩
      exact ⟨elem, ih rest⟩

-- this is also an unused theorem.
theorem span_ordering : ∀(A B : Set V), A ⊆ B → span A ⊆ span B := by
  intro A B hypsubset
  unfold span
  simp
  intros v ls hypeq hypall
  constructor
  · constructor
    · apply hypeq
    · exact allin_ordering ls A B hypsubset hypall

theorem zero_in_span : ∀(S : Set V), 0 ∈ span S := by
  intro S
  unfold span
  let ex : List (ℝ × V) := List.nil
  dsimp
  -- ⊢ ∃ C, sumUp C = 0 ∧ allin C S
  fconstructor -- constructor for ∃
  · exact ex -- example empty list
  · -- unfold definitions to show it's a linear combo for 0
    unfold ex sumUp allin
    simp

theorem span_mul_closed : ∀(S : Set V) (r : ℝ)
    (u : {v : V // v ∈ span S}), r • u.val ∈ span S := by
  intros S r u
  unfold span; dsimp
  -- ⊢ ∃ C, sumUp C = r • ↑u ∧ allin C S

  let pu := u.prop
  unfold span at pu; dsimp at pu
  -- pu : ∃ C, sumUp C = ↑u ∧ allin C S
  --       ⟨lu, sumUp lu = ↑u ∧ allin lu S⟩
  rcases pu with ⟨lu, pu⟩
  let ex: List (ℝ × V) := lsmul r lu
  have sum_val : sumUp ex = r • u := by
    unfold ex
    -- ⊢ sumUp (lsmul r lu) = r • ↑u
    rw [sumup_lin, pu.left]
  have allin_proof : allin ex S := by
    unfold ex
    -- ⊢ allin (lsmul r lu) S
    exact (lsmul_allin lu S r).mp pu.right
  constructor -- for ∃
  · exact ⟨sum_val, allin_proof⟩

theorem span_add_closed : ∀(S : Set V) (a b : {v : V // v ∈ span S}),
    a.val + b.val ∈ span S := by
  intros S a b
  let pa := a.prop -- a ∈ span S
  let pb := b.prop -- b ∈ span S
  unfold span at pa; dsimp at pa
  unfold span at pb; dsimp at pb

  -- pa : ∃ C, sumUp C = ↑a ∧ allin C S
  --      ⟨xa, ⟨sumUp xa = ↑a, allin xa S⟩⟩
  rcases pa with ⟨xa, ⟨ha_sum, ha_sub⟩⟩
  -- pb : ∃ C, sumUp C = ↑b ∧ allin C S
  --      ⟨xb, ⟨sumUp xb = ↑b, allin xb S⟩⟩
  rcases pb with ⟨xb, ⟨hb_sum, hb_sub⟩⟩
  have sumeq: sumUp (xa ++ xb) = a.val + b.val := by
    calc
      _ = sumUp xa + sumUp xb := by exact sumup_concat xa xb
      _ = a.val + b.val := by rw [ha_sum, hb_sum]
  have ain: allin (xa ++ xb) S := by
    exact (allin_concat xa xb S).mp ⟨ha_sub, hb_sub⟩

  unfold span
  dsimp
  -- ⊢ ∃ C, sumUp C = ↑a + ↑b ∧ allin C S
  constructor
  · exact ⟨sumeq, ain⟩

-- define the zero vector on the span subspace
instance (S : Set V) : Zero {v : V // v ∈ span S} where
  zero := Subtype.mk 0 <| zero_in_span S

-- vector negation
instance (S : Set V) : Neg {v : V // v ∈ span S} where
  neg (a : {v : V // v ∈ span S}) :=
    Subtype.mk (-a.val) <| by
      have sub : -a.val = (-1 : ℝ) • a.val := by simp
      rw [sub]
      exact span_mul_closed S (-1) a

-- vector addition
instance (S : Set V) : Add {v : V // v ∈ span S} where
  add (a b : {v: V // v ∈ span S}) :=
    Subtype.mk (a.val + b.val) <| span_add_closed S a b

@[simp]
theorem subtype_ext_add : ∀(S: Set V) (a b : {v : V // v ∈ span S}),
    (a + b).val = a.val + b.val := by intros S a b; rfl

@[simp]
theorem subtype_ext_neg : ∀(S : Set V) (a : {v : V // v ∈ span S}),
  (-a).val = -(a.val) := by intros S a; rfl

@[simp]
theorem subtype_ext_zero : ∀(S : Set V),
  (0 : {v : V // v ∈ span S}).val = 0 := by intro S; rfl

-- vector addition is an abelian group
instance (S : Set V) : AddCommGroup {v: V // v ∈ span S} where
  add_assoc := by intro a b c; ext; simp; rw [add_assoc]
  zero_add := by intro z; ext; simp;
  add_zero := by intro a; ext; simp;
  add_comm := by intro a b; ext; simp; rw [add_comm]
  neg_add_cancel := by intro a; ext; simp;
  nsmul := nsmulRec
  zsmul := zsmulRec

-- scalar multiplication
instance (S : Set V) : SMul ℝ {v : V // v ∈ span S} where
  smul (r : ℝ) (u : {v : V // v ∈ span S}) :=
    Subtype.mk (r • u.val) <| span_mul_closed S r u

@[simp]
theorem subtype_ext_smul : ∀(S: Set V) (r: ℝ)
    (u : {v : V // v ∈ span S}), (r • u).val = r • u.val := by
  intros S a b; rfl

-- we have the module axioms: spans are subspaces!
instance (S : Set V) : Module ℝ {v : V // v ∈ span S} where
  one_smul := by intro b; ext; simp
  mul_smul := by intros a b c; ext; simp; rw [mul_smul]
  smul_zero := by intros a; ext; simp;
  smul_add := by intros a b c; ext; simp
  add_smul := by intros a b c; ext; simp; rw [add_smul]
  zero_smul := by intros a; ext; simp

def nallz (ls : List (ℝ × V)) : Prop := match ls with
| List.cons ⟨r, _⟩ rest => r ≠ 0 ∨ nallz rest
| List.nil => False

def linIndep (S : Set V) : Prop :=
    ¬∃(C : List (ℝ × V)), sumUp C = 0 ∧ nallz C ∧ allin C S

structure Basis (S : Set V) where
  spansFull : ∀(v: V), v ∈ span S
  indep : linIndep S

/- im convinced this is impossible
theorem dim_invariant : ∀(B₁ B₂ : Set V),
    (h₁ : Basis B₁) → (h₂ : Basis B₂) →
      ∃(f : {v : V // v ∈ B₁} → {v : V // v ∈ B₂}), Bijective f := by
  intros B₁ B₂ h₁ h₂
  let contra:
      (¬∃(f : {v : V // v ∈ B₁} → {v : V // v ∈ B₂}), Bijective f)
        → False := by
    intro hyp
    simp at hyp
    sorry
  exact Classical.byContradiction contra
-/

end Vect

    \end{lstlisting}


\end{document}
